\documentclass[11pt,a4paper]{article}
\usepackage[margin=24mm]{geometry}
\usepackage{amsmath,amssymb,bm,booktabs,longtable,array,hyperref,xcolor}
\hypersetup{colorlinks=true,linkcolor=teal,urlcolor=teal,citecolor=teal,pdftitle={Two Weyssenhoff fluids: spherical equations and bounce audit}}
\newcommand{\dd}{\mathrm d}
\newcommand{\Dt}{D_t}
\newcommand{\Dr}{D_r}
\newcommand{\kap}{\kappa}
\newcommand{\eps}{\epsilon}
\newcommand{\Tr}{\mathcal T}
\setlength{\emergencystretch}{3em}
\title{Two Weyssenhoff fluids in general spherical Einstein--Cartan collapse\\\large Stage 2: equations, bounce conditions and mathematical viability}
\author{EC Collapse Laboratory research record\\\small\texttt{pdamkier-del/skovlab}}
\date{12 November 2026}
\begin{document}
\maketitle
\begin{abstract}
Continuing the audit at commit \texttt{3efc961}, we compare a conditionally local averaged Hashemi-type Weyssenhoff fluid with the full ideal field formulation of Obukhov and Korotky reproduced by Br\'echet, Hobson and Lasenby. The lapse, radial metric and areal radius remain general. Proper spherical symmetry and Frenkel leave one signed coherent radial spin density in the second model. We derive the cancellation of its linear derivative stress only after symmetry reduction. Both selected closures have the same acoustic pair, with additional advected composition or polarization modes. Conserved-spin ordinary dust is nonhyperbolic on its low-spin branch; enthalpy and zero-sound-speed surfaces are degenerate. We derive distinct shell and volume turning conditions, local curvature checks, analytical constrained initial patches and necessary vacuum-interface conditions. These do not establish a causal regular finite-cloud bounce. They are conditional fluid conclusions, not a no-go theorem for minimal Einstein--Cartan gravity. No numerical evolution, plots, simulator or UI changes are made.
\end{abstract}
\tableofcontents
\section{Selected formulation, action and evidence status}
\paragraph{Scope.} Model A here is the conditionally local, statistically unpolarized Hashemi-type fluid in general spherical geometry. Its homogeneous FLRW reference module is preserved unchanged. Model B is the ideal, neutral, minimally coupled, convective Frenkel fluid in \cite{BHL,OK}, with all fields inheriting proper spherical rotations. It does not mean a new dissipative spin theory, arbitrary spin moments, or a merely spherical metric with nonsymmetric hidden sources. We distinguish \emph{literature-established field equations}, \emph{derived consequences}, \emph{conditional material assumptions}, and \emph{unresolved physics}. The last three categories are not promoted to microscopic predictions by a symbolic check.

\subsection{Conventions and bulk variation}
Use $c=1$, $\kap=8\pi G$, signature $(-+++)$, weight-$1/2$ brackets, and the derivative index first:
\begin{align}
 \widetilde\nabla_\mu V^\rho&=\partial_\mu V^\rho+\widetilde\Gamma^\rho{}_{\mu\nu}V^\nu,&
 \Tr^\rho{}_{\mu\nu}&=2\widetilde\Gamma^\rho{}_{[\mu\nu]},\\
 K^\rho{}_{\mu\nu}&=\widetilde\Gamma^\rho{}_{\mu\nu}-\Gamma^\rho{}_{\mu\nu},&
 K_{abc}&=\tfrac12(\Tr_{abc}-\Tr_{bca}+\Tr_{cab}).
\end{align}
Here $K_{abc}=-K_{cba}$, untilded $\nabla$ is Levi--Civita, and
\begin{equation}
 \widetilde R^\rho{}_{\sigma\mu\nu}=\partial_\mu\widetilde\Gamma^\rho{}_{\nu\sigma}-\partial_\nu\widetilde\Gamma^\rho{}_{\mu\sigma}
 +\widetilde\Gamma^\rho{}_{\mu\lambda}\widetilde\Gamma^\lambda{}_{\nu\sigma}
 -\widetilde\Gamma^\rho{}_{\nu\lambda}\widetilde\Gamma^\lambda{}_{\mu\sigma},\quad
 \widetilde R_{\sigma\nu}=\widetilde R^\rho{}_{\sigma\rho\nu}.
\end{equation}
The action has no Holst term, cosmological constant, nonminimal spin coupling or propagating torsion term:
\begin{equation}
 I=\frac1{2\kap}\int\sqrt{-g}\,\widetilde R\,\dd^4x+I_m[e,\widetilde\omega,\psi].
\end{equation}
With $\epsilon_{0123}=+1$, $\eta=e^0\wedge e^1\wedge e^2\wedge e^3$, define sources by
\begin{equation}
 \delta I_m=\int\delta e^a\wedge\Sigma_a-\tfrac12\int\delta\widetilde\omega^{ab}\wedge\tau_{ab}+\int E_\psi\delta\psi.
\end{equation}
Varying the gravitational 4-form $(4\kap)^{-1}\epsilon_{abcd}e^a\wedge e^b\wedge\widetilde{\mathcal R}^{cd}$, using $\delta\mathcal R=\widetilde D\delta\omega$ and one integration by parts for compactly supported variations, gives
\begin{equation}
 \epsilon_{abcd}e^b\wedge\widetilde{\mathcal R}^{cd}=-2\kap\Sigma_a,
 \qquad \epsilon_{abcd}\Tr^a\wedge e^b=\kap\tau_{cd},\qquad E_\psi=0.
\end{equation}
Local Lorentz invariance yields $\widetilde D\tau_{ab}+2e_{[a}\wedge\Sigma_{b]}=0$. These are bulk variational equations; interface terms still require a boundary problem. The convective source $\tau^\lambda{}_{\mu\nu}=u^\lambda s_{\mu\nu}$, $s_{\mu\nu}u^\nu=0$, gives trace-free algebraic torsion $\Tr^\lambda{}_{\mu\nu}=\kap u^\lambda s_{\mu\nu}$. There is no independent torsion characteristic cone. The Stage 1 rest-frame auxiliary-variable calculation checks the quadratic coefficient: $\Delta\widetilde R=-\kap^2s^2/2$ and the on-shell sum of gravitational and linear material contortion terms is $+\kap s^2/4$.

This variation alone does not specify $I_m$, its thermodynamics, or its director constraints. We identify the established full effective equations of \cite{BHL}, rather than invent the original constrained action of \cite{OK}. An independent reconstruction of that full action remains an open reproducibility requirement; the calculations below have field-formulation status.

\subsection{Full effective field equations before symmetry reduction}
Retain the BHL spin variable $b_{\mu\nu}$ and $b^2=\tfrac12b_{\mu\nu}b^{\mu\nu}$. BHL use signature $+---$, derivative index last, half-torsion $Q_B$, and $\widetilde\Gamma_B=\Gamma_B-K_B$. For the same geometric connection our dictionary is
\begin{equation}
 \Tr=-2Q_B,\quad K^\rho{}_{\mu\nu}=-(K_B)^\rho{}_{\nu\mu},\quad s_{\mu\nu}=-2b_{\mu\nu},\quad
 q=\kap b^2=\kap s^2/4.
\end{equation}
Appendix B of \cite{BHL} gives $g_B=-g$, $u_B^\mu=u^\mu$, $(u_B)_\mu=-u_\mu$, $(b_B)_{\mu\nu}=b_{\mu\nu}$. Transforming their Eq. (24) therefore yields
\begin{align}
 G_{\mu\nu}&=\kap T^{\rm eff}_{\mu\nu},\label{eq:fullEFE}\\
 T^{\rm eff}_{\mu\nu}&=(\rho+p-2q)u_\mu u_\nu+(p-q)g_{\mu\nu}+\mathcal D_{\mu\nu},\label{eq:fullT}\\
 \mathcal D_{\mu\nu}&=-2(g^{\alpha\beta}-u^\alpha u^\beta)\nabla_\alpha[u_{(\mu}b_{\nu)\beta}].\label{eq:D}
\end{align}
The factor $g^{\alpha\beta}-u^\alpha u^\beta$ in this transformed expression is not the spatial projector. The spin current in BHL units is $u^\lambda b_{\mu\nu}$; the operational source used in the bulk variation is $-2u^\lambda b_{\mu\nu}$. Hashemi's existing input subtracts $q=\kap S_H^2/4$; aligning scalar corrections gives $S_H^2=4b^2$ only for the same coupling/source assumptions. It does not identify coherent spin with ensemble variance.

Let $h_\mu{}^\nu=\delta_\mu^\nu+u_\mu u^\nu$ and $a_\mu=u^\lambda\nabla_\lambda u_\mu$. The ideal spin equation and Frenkel constraint are
\begin{align}
 h_\mu{}^\alpha h_\nu{}^\beta\nabla_\lambda(u^\lambda b_{\alpha\beta})&=0,
 &b_{\mu\nu}u^\nu&=0,\label{eq:spinproject}\\
 \Dt b_{\mu\nu}+\Theta b_{\mu\nu}&=-2u_{[\mu}b_{\nu]\lambda}a^\lambda,
 &\Dt b^2&=-2\Theta b^2.\label{eq:spintransport}
\end{align}
The second line follows by differentiating Frenkel and projecting the first. Its contraction proves the scalar law, with no relaxation or torque term. In addition impose $\nabla_\mu T^{\mu\nu}_{\rm eff}=0$; do not replace this by a guessed bare canonical conservation law. Particle conservation $\nabla_\mu(nu^\mu)=0$ and an ordinary material EOS are explicitly stated ideal material assumptions when $n$ is used below.

\section{Model A: statistically unpolarized matter in GSM}
\subsection{Microscopic torsion versus an averaged effective source}
Use Hashemi's spin normalization $S^H_{\mu\nu}$ with
\begin{equation}
 \langle S^H_{\mu\nu}\rangle=0,\qquad
 S_H^2=\tfrac12\langle S^H_{\mu\nu}S_H^{\mu\nu}\rangle,\qquad
 q=\kap S_H^2/4.
\end{equation}
For the aligned ideal minimal coupling, the microscopic source in our operational convention may be denoted $s=S^H$, satisfying Frenkel. Cartan and contortion then have the explicit forms
\begin{equation}
 \Tr^\rho{}_{\mu\nu}=\kap u^\rho s_{\mu\nu},\qquad
 K_{abc}=\tfrac\kap2(u_a s_{bc}-u_b s_{ca}+u_c s_{ab}).
\end{equation}
In a common-velocity ensemble, $\langle\Tr\rangle=\langle K\rangle=0$, but quadratic products have nonzero average. Eliminate microscopic torsion \emph{before} averaging; setting torsion identically zero first gives a different theory. The effective averaged metric equations do not specify a unique resolved torsionful connection or its full curvature statistics. Those require microscopic/correlation information beyond $q$.

For Model B the same formula uses the coherent operational spin $s=-2b$, not the same numerical tensor as Hashemi's named spin. Thus $q_B=\kap b^2$ and $q_A=\kap S_H^2/4$ align scalar corrections under the stated source assumptions; the coherent mean is not the ensemble variance.

\subsection{Additional local closure assumptions}
The homogeneous expressions in \cite{Hashemi} do not imply their arbitrary-gradient extension. We explicitly assume: a separation of averaging and resolved gradient scales; a shared resolved congruence; locally rotationally isotropic microscopic spin statistics; negligible unresolved velocity--spin, acceleration--spin, shear--spin and connection correlations; commuting averaging and covariant derivatives at retained order; negligible metric averaging backreaction; conserved particles and spin per particle with no torque, exchange or production. These are material approximations, not consequences of the spherical metric. Under them the mean linear-spin derivative source vanishes and
\begin{equation}\label{eq:Aconstitutive}
 T^A_{\mu\nu}=(\rho+p-2q)u_\mu u_\nu+(p-q)g_{\mu\nu},\quad
 \eps_A=\rho-q,\quad P_A=p-q,\quad Q_A=0.
\end{equation}
The local isotropic covariance in Hashemi units is
\begin{equation}
 \langle S^H_{ij}S^H_{kl}\rangle
 =\tfrac13S_H^2(\delta_{ik}\delta_{jl}-\delta_{il}\delta_{jk}).
\end{equation}
It permits nonzero smooth variance at the center. It is not the coherent radial field of Model B.

A checkable reduced action uses a conserved-current perfect-fluid variational principle of Brown type \cite{Brown}, with current density $J^\mu$ and advected composition $\chi$:
\begin{equation}
 I_{\rm eff}=\int\left[-\sqrt{-g}\,\eps(n,\chi)
 +J^\mu(\partial_\mu\varphi+\lambda\partial_\mu\chi+
 \lambda_A\partial_\mu\alpha^A)\right]\dd^4x,\quad
 n=\frac{\sqrt{-g_{\mu\nu}J^\mu J^\nu}}{\sqrt{-g}}.
\end{equation}
Choose, as the declared effective reference closure,
\begin{equation}
 \eps(n,\chi)=C n^{1+w}-\chi n^2,\qquad
 \rho=C n^{1+w},\quad q=\chi n^2,\quad C>0,\quad\chi\geq0.
\end{equation}
Variation at fixed current gives $P=n\eps_n-\eps=w\rho-q$, particle conservation and $\Dt\chi=0$. It follows that
\begin{equation}\label{eq:Atransport}
 \Dt\rho=-(1+w)\rho\Theta,\quad\Dt q=-2q\Theta,\quad
 \partial_t(BR^2n)=0,\quad\Dt(q/n^2)=0.
\end{equation}
This action verifies the conditional \emph{effective metric-fluid} theory. It is not a reconstructed microscopic Weyssenhoff matter action or proof of the averaging assumptions.

\subsection{Independent spherical equations and comparison}
Insert Eq.~\eqref{eq:Aconstitutive} in all mixed metric components in the geometry section. The equations for $R,\Phi,\Lambda$, mass and constraints are the same perfect-source equations as for the reduced coherent Model B, using $q_A$ instead of $\kap b^2$. The independent matter variables are $(\rho,q)$ or $(n,\chi)$, not both pairs. Initial profiles after $R_i=r$ are $(U_i,\rho_i,q_i)$; pressure, effective density and constrained geometry follow. In particular
\begin{equation}
 h=(1+w)\rho-2q,\quad
 \Phi'=-\frac{w\rho'-q'}h,\quad
 H_r=\Dr U/\Gamma,\quad\Dt R=U
\end{equation}
on a nondegenerate monotone patch. These equations retain pressure acceleration and shear. The undivided conservation equation $\Dr P+h a_r=0$ remains valid when the divided expression fails.

There is a complete averaged-fluid radial pencil in $(\rho,q,v)$:
\begin{equation}
 A_A^0=\operatorname{diag}(1,1,h),\qquad
 A_A^1=\begin{pmatrix}0&0&(1+w)\rho\\0&0&2q\\w&-1&0\end{pmatrix}.
\end{equation}
Its determinant is $\lambda[w(1+w)\rho-2q-h\lambda^2]$. Alternatively $(n,\chi,v)$ gives Eq.~\eqref{eq:pencil} with $P_\chi=-n^2$. Thus the characteristic classification and enthalpy problem apply to Model A independently of any identification with coherent spin. At zero variance the coherent variable change $q=\kap b^2$ is not invertible; the separate pencils are the correct way to compare the models there.

\begin{longtable}{p{.22\textwidth}p{.33\textwidth}p{.34\textwidth}}
\toprule Quantity & Model A: averaged GSM & Model B: full coherent GSM\\\midrule\endhead
Spin content & Zero mean; isotropic nonzero variance. & Signed radial axial mean $b$ under proper rotations.\\
Quadratic correction & $\kap S_H^2/4$ & $\kap b^2$ in BHL units; not a factor to copy blindly.\\
Linear-spin terms & Vanish only under the stated averaging/commutation assumptions. & Retained first, then vanish by exact radial symmetry.\\
Flux/anisotropy & Zero under the selected moment truncation. & Zero in the exact ideal radial branch, despite shear and gradients.\\
Independent spin transport & $\Dt q=-2q\Theta$ & $\Dt b=-b\Theta$, hence $\Dt q=-2q\Theta$.\\
Center & Even variance may be nonzero. & $b=O(r)$, so coherent $q=O(r^2)$.\\
Closure status & Conditional statistical/material approximation. & Conditional selected full ideal field formulation plus EOS; not a generic correlation hierarchy.\\
Characteristics & Contact/composition mode and acoustic pair. & Contact/polarization mode and same pair away from degeneracies.\\
Full-connection boundary & Microscopic spin statistics/edge layer not determined by mean-zero torsion alone. & Restricted contortion-jump test available; general spinful layers unresolved.\\\bottomrule
\end{longtable}

\section{Spin gradients, shear, acceleration and spherical symmetry}
\subsection{Retaining the linear terms until after reduction}
For a hypersurface-orthogonal congruence write
\begin{equation}
 \nabla_\alpha u_\mu=-u_\alpha a_\mu+\tfrac13\Theta h_{\alpha\mu}+\sigma_{\alpha\mu},\qquad
 \mathsf D_\alpha b_{\mu\nu}=h_\alpha{}^\lambda h_\mu{}^\rho h_\nu{}^\sigma\nabla_\lambda b_{\rho\sigma}.
\end{equation}
Expanding the product in Eq.~\eqref{eq:D} and using the differentiated Frenkel condition gives
\begin{align}
 J_\mu&=\mathsf D^\lambda b_{\mu\lambda}+2b_{\mu\lambda}a^\lambda,\\
 \mathcal D_{\mu\nu}&=-2u_{(\mu}J_{\nu)}-2\sigma_{(\mu}{}^\lambda b_{\nu)\lambda}.\label{eq:decomp}
\end{align}
These equations are restricted to zero vorticity, as required by the comoving spherical geometry, not a general rotating-fluid decomposition. In a local orthonormal rest frame with diagonal expansion rates $H_i$,
\begin{equation}
 \mathcal D_{\hat0\hat0}=0,\quad
 \mathcal D_{\hat0\hat i}=\mathsf D^j b_{ij}+2b_{ij}a^j,\quad
 \mathcal D_{\hat i\hat j}=(H_i-H_j)b_{ij}\quad(i\ne j).
\end{equation}
For example, $b_{12}\ne0$ with $H_1\ne H_2$ has nonzero shear stress. The cancellation below cannot be obtained by deleting all spin--shear terms at the outset. In our convention an effective energy flux is $Q_i=-\mathcal D_{\hat0\hat i}=-J_i$.

\subsection{Exact coherent spherical branch}
Set $A=e^\Phi$, $B=e^\Lambda$, $u=A^{-1}\partial_t$. An invariant two-form under the three proper rotations has only the two possible blocks
\begin{equation}
 b=f(t,r)\,\dd t\wedge\dd r+g(t,r)\sin\theta\,\dd\theta\wedge\dd\phi.
\end{equation}
The Frenkel condition kills $f$. Parametrize the remaining block as
\begin{equation}
 b_{\theta\phi}=R^2\sin\theta\,b(t,r),\quad b_{\phi\theta}=-b_{\theta\phi},\quad b^2=b(t,r)^2.\label{eq:radial}
\end{equation}
The signed scalar on the right is an axial radial polarization, not an unpolarized spin magnitude. $SO(3)$ permits it. Requiring also parity invariance forces its mean to zero: an axial vector transforms differently from the polar radial direction. At a regular center, smooth Cartesian components require the axial vector $\bm b=F(t,r^2)\bm x$, hence $b=O(r)$. A nonzero constant radial magnitude is not smooth there.

On each spatial slice, $b_{ij}$ is dual to the radial vector $b n^i$. Its spatial divergence is the curl of that radial vector, which vanishes for arbitrary radial $b(t,r)$; this remains true for the curved spherical spatial metric. The acceleration is radial, so $b_{ij}a^j=0$. Finally the angular expansion eigenvalues coincide, $H_\theta=H_\phi=H_\perp$, while the only spin plane is the angular plane. Each term in Eq.~\eqref{eq:decomp} therefore vanishes separately:
\begin{equation}
 \mathcal D_{\mu\nu}=0,\quad Q=0,\quad P_r=P_\perp=P=p-q,\quad \eps=\rho-q,\quad q=\kap b^2.\label{eq:radsource}
\end{equation}
This derivation permits arbitrary $b'$, $\Phi'$, inhomogeneous density and radial shear. It is not a homogeneous approximation. Direct four-dimensional covariant evaluation and separate spatial-divergence, acceleration and shear tests independently check it.

The angular component of Eq.~\eqref{eq:spinproject} now becomes
\begin{equation}
 \partial_t b+(\partial_t\Lambda+2\partial_tR/R)b=0,
 \qquad \partial_t(BR^2b)=0.\label{eq:radtransport}
\end{equation}
Particle conservation gives $\partial_t(BR^2n)=0$, so the signed spin per particle $\beta=b/n$ is advected, $\Dt\beta=0$. Extra coherent transverse spin modes are forbidden in this exact symmetry class. Nonspherical perturbations of the same background require a different, multidimensional principal analysis; our radial calculation makes no assertion of stability against them.

\subsection{Coherent means, variance and correlations are different data}
An ensemble with common resolved velocity and metric may have mean axial spin $m_i=b n_i$ and raw second moment
\begin{equation}
 C_{ij}=C_\parallel n_i n_j+C_\perp(h_{ij}-n_i n_j),\quad
 C_\perp\geq0,\quad C_\parallel\geq b^2.
\end{equation}
The total squared spin is $\langle b^2\rangle=C_\parallel+2C_\perp$; the central variance is this trace minus $b^2$. Smooth isotropic statistics at the center have $C_\parallel=C_\perp$ there and may have nonzero variance, unlike the coherent radial mean.

Under \emph{common congruence, commuting derivatives and averaging, no unresolved metric/velocity correlations, and the same ideal transport and ordinary EOS}, Fermi--Walker components obey $\Dt C_{ij}=-2\Theta C_{ij}$. Averaging Eq.~\eqref{eq:fullT} retains only $q=\kap\langle b^2\rangle$ and the linear derivative of the mean, which still cancels for a radial mean. Even retaining $C_\parallel,C_\perp$ under these assumptions introduces only advected composition modes; the BHL algebraic quadratic term depends on their trace and does not create an anisotropic pressure. Isotropy of the second moment is not needed for this specific algebraic statement. A spin-dependent ordinary internal energy would instead require a new, specified material EOS.

If velocities, shear, acceleration or geometry fluctuate, averages such as $\langle\delta u\,\delta b\rangle$, $\langle\delta\sigma\,\delta b\rangle$ and connection--spin terms need not vanish. Averaging the convective spin equation then produces flux and higher moments, not a transport law for their trace alone. Spherical symmetry allows scalar radial flux and pressure anisotropy, but supplies no evolution or susceptibility for them. An exact source decomposition with symbols $Q,P_r,P_\perp$ is not closure. No characteristic determinant can be assigned to this unclosed hierarchy. We have not established that all possible microscopic theories fail; we have established that naming these correlations does not repair the selected ideal system.

\subsection{Relation to FLRW}
Full-field FLRW invariance excludes any nonzero coherent convective Frenkel two-form, since its spatial dual selects a direction. A statistically isotropic mean-zero variance supports Model A under its documented averaging assumptions. Stage 1 also checked a weaker, flat FLRW \emph{metric} with homogeneous Cartesian aligned spin, $b_{xy}=a^2\beta_0(t)$, $\dot\beta_0=-3H\beta_0$: its effective derivative source vanishes although its spin/torsion is not isotropic. That example does not evade the strong-symmetry restriction and is not a radial center-smooth full-field FLRW limit of Eq.~\eqref{eq:radial}.

\section{General spherical gravitational and matter equations}
\subsection{Documented metric derivation and mixed components}
Write $g=h_{ab}\dd x^a\dd x^b+R^2\gamma_{ij}\dd y^i\dd y^j$, with $h=\operatorname{diag}(-A^2,B^2)$ and $\gamma$ the unit sphere. The mixed Christoffels are $\Gamma^a{}_{ij}=-RR^{,a}\gamma_{ij}$ and $\Gamma^i{}_{aj}=R_{,a}\delta^i_j/R$. Defining $X=h^{ab}R_{,a}R_{,b}$ and $H_{ab}=\nabla^{(h)}_a\nabla^{(h)}_bR$, insertion in the Ricci definition gives
\begin{align}
 R^{(4)}_{ab}&=\tfrac12R_hh_{ab}-2H_{ab}/R,&
 R^{(4)}_{ij}&=(1-X-R\Box_hR)\gamma_{ij},\\
 G_{ab}&=-2H_{ab}/R+[2\Box_hR/R+(X-1)/R^2]h_{ab},&
 G^\theta{}_\theta&=\Box_hR/R-R_h/2.
\end{align}
Using the two-dimensional Christoffels to expand these Hessians yields the requested mixed components (dots/primes are coordinate derivatives):
\begin{align}
 G^t{}_t={}&-A^{-2}\left(2\dot R\dot\Lambda/R+\dot R^2/R^2\right)
 +B^{-2}\left(2R''/R+R'^2/R^2-2\Lambda'R'/R\right)-R^{-2},\label{eq:Gtt}\\
 G^r{}_r={}&-A^{-2}\left(2\ddot R/R+\dot R^2/R^2-2\dot\Phi\dot R/R\right)
 +B^{-2}\left(R'^2/R^2+2\Phi'R'/R\right)-R^{-2},\\
 G^t{}_r={}&\frac{2}{A^2R}(\dot R'-\Phi'\dot R-\dot\Lambda R'),\\
 G^\theta{}_\theta={}&-A^{-2}\left[\ddot\Lambda+\dot\Lambda^2-\dot\Phi\dot\Lambda+
 (\ddot R+(\dot\Lambda-\dot\Phi)\dot R)/R\right]\nonumber\\
 &+B^{-2}\left[\Phi''+\Phi'^2-\Phi'\Lambda'+(R''+(\Phi'-\Lambda')R')/R\right].\label{eq:Gangle}
\end{align}
For the coherent radial source their right-hand sides are respectively $-\kap\eps$, $\kap P$, $0$, $\kap P$. For a general outward energy flux $Q$ with $T_{\hat0\hat1}=-Q$, instead $G^t{}_r=\kap(B/A)Q$. Raising the first time index changes the sign; the covariant $G_{tr}$ must not be mistaken for $G^t{}_r$. The CAS independently constructs the full four-dimensional metric connection and checks this conversion as well as Stage 1's components.

\subsection{Undivided conservation and spin equations}
Let $\Dt=A^{-1}\partial_t$, $\Dr=B^{-1}\partial_r$, $U=\Dt R$, $\Gamma=\Dr R$, $H_r=\Dt\Lambda$, $H_\perp=U/R$ and $a_r=\Dr\Phi$. Then
\begin{equation}
 \Theta=H_r+2H_\perp,\qquad \sigma^2=\tfrac13(H_r-H_\perp)^2.
\end{equation}
Before adopting a perfect source, the spherical Bianchi/conservation projections read
\begin{align}
 \Dt\eps+(\eps+P_r)H_r+2(\eps+P_\perp)H_\perp
 +\Dr Q+2Q(\Gamma/R+a_r)&=0,\label{eq:generalenergy}\\
 \Dt Q+2Q(H_r+H_\perp)+\Dr P_r+(\eps+P_r)a_r
 +2\Gamma(P_r-P_\perp)/R&=0.\label{eq:generalmomentum}
\end{align}
For Eq.~\eqref{eq:radsource} these reduce, without division by enthalpy, to
\begin{equation}
 \Dt\eps+h\Theta=0,\quad \Dr P+h a_r=0,\quad h=\eps+P=\rho+p-2q.
\end{equation}
Together with $\Dt q=-2q\Theta$ from the full spin equation, they imply
\begin{equation}
 \Dt\rho=-(\rho+p)\Theta,\quad \Dt b=-b\Theta,\quad
 \Dt n=-n\Theta,\quad \Dt(b/n)=0.\label{eq:material}
\end{equation}
The $n$ equation is used only with a declared conserved-particle material interpretation; for fixed $p=w\rho$ the independent pair $\rho,b$ already suffices, and $n$ must not be counted as an additional arbitrary thermodynamic profile. Equivalently use $\rho=C n^{1+w}$ and advected $\beta=b/n$, with fixed $C>0$ in the isentropic case. An independent entropy variable requires its EOS and ideal advection equation. Effective conservation by itself would not evolve an arbitrary separate spin variance.

\subsection{Constraints, mass and candidate evolution system}
Define geometric Misner--Sharp mass $M=Gm_{\rm MS}=R(1+U^2-\Gamma^2)/2$. Differentiating it and using the metric equations gives, for the radial coherent branch,
\begin{align}
 \Dr M&=\tfrac\kap2\eps R^2\Gamma,&\Dt M&=-\tfrac\kap2 P R^2U,\\
 \Dr U-\Gamma H_r&=0,&\Dt\Gamma&=U a_r,\\
 \Dt R&=U,&\Dt U&=\Gamma a_r-M/R^2-\tfrac\kap2 PR.\label{eq:gravityevolution}
\end{align}
The ADM Hamiltonian constraint, with $\mathsf K^i{}_j=-\operatorname{diag}(H_r,H_\perp,H_\perp)$, is
\begin{equation}
 {}^{(3)}R+4H_rH_\perp+2H_\perp^2=2\kap\eps,\quad
 {}^{(3)}R=2(1-\Gamma^2)/R^2-4\Dr\Gamma/R.
\end{equation}
For retained flux the exact identities become $\Dr M=\kap R^2(\eps\Gamma+QU)/2$, $\Dt M=-\kap R^2(P_rU+Q\Gamma)/2$, $\Dr U-\Gamma H_r=\kap RQ/2$ and $\Dt\Gamma=Ua_r+\kap RQ/2$. They are verified geometry identities, not a full-correlation constitutive theory.

On a monotone material chart with $\Gamma\ne0$ and $h\ne0$, a possible constrained state is $(R,U,\rho,b)$, with
\begin{align}
 \Theta&=\Dr U/\Gamma+2U/R,& a_r&=-\Dr P/h,\\
 M'&=\tfrac\kap2\eps R^2R',&\Gamma^2&=1+U^2-2M/R,\\
 B&=R'/\Gamma,&\Phi'&=-P'/h.
\end{align}
Evolve $R,U,\rho,b$ using Eqs.~\eqref{eq:material}, \eqref{eq:gravityevolution}; for Model A replace $b$ by $q$ and its variance transport. In coordinate form $\partial_tR=AU$ and $\partial_t\Lambda=A\,\Dr U/\Gamma$, while $\partial_r\Phi=-\partial_rP/h$ fixes the lapse up to one time normalization. Thus no independent material evolution law for $\Phi$ is missing: it is a constrained gauge variable in this split. Integrate the radial mass and lapse constraints, and monitor the unused mass-time, angular and ADM equations. This is a formal continuum split, not a tested numerical algorithm. Do not independently overwrite all redundant fields. Neither the $h=0$ nor $\Gamma=0$ surfaces can be crossed by this divided scheme; the undivided equations are the correct starting point for any separate analysis there. Bianchi identities establish continuum constraint compatibility conditional on matter evolution, not discrete stability or a regular free boundary.

\subsection{Reference limits}
With $b=0$, the equations reduce to spherical GR Einstein--Euler and pre-caustic LTB dust. With $R=ar$, $B=a/\sqrt{1-kr^2}$, $A=1$, the metric equations become
\begin{equation}
 H^2+k/a^2=\kap(\rho-q)/3,\quad
 \ddot a/a=-\kap[(1+3w)\rho-4q]/6,\quad
 \dot\rho=-3(1+w)H\rho,\quad\dot q=-6Hq.
\end{equation}
For nonzero coherent radial spin this is only a metric-source limit; a regular full-field FLRW spin configuration does not result. Model A uses ensemble variance instead, $q=\kap S_H^2/4$, with its existing implementation and coefficients preserved.

\section{Radial characteristic and high-density analysis}
\subsection{Full reduced principal symbol, including polarization}
This is a local, frozen-coefficient, high-frequency test at a noncentral regular event. The ordinary causal-fluid reference restricts $0\leq w\leq1$; the algebraic pencils remain valid for other real $w$, but no complete physical classification outside this reference range is claimed. Use a locally inertial fluid rest frame $(\tau,x)$ with $x$ radial. Curvature, spherical $1/R$ terms, background gradients and nonderivative gravity terms are lower order. For radial perturbations the spin remains in the angular plane; Frenkel is preserved under radial boosts. Choose the isentropic ordinary EOS $\rho=Cn^{1+w}$, conserved $n>0$, and independent signed $\beta=b/n$. Thus
\begin{equation}
 \eps=Cn^{1+w}-\kap\beta^2n^2,\quad
 P=wCn^{1+w}-\kap\beta^2n^2,\quad h=n\eps_{,n}=(1+w)\rho-2q.
\end{equation}
The equality $P=n\eps_{,n}-\eps$ provides a thermodynamic cross-check on this reduced source, not a replacement microscopic Weyssenhoff action. Linearized particle, polarization and momentum equations give
\begin{equation}
 A^0\partial_\tau\delta W+A^1\partial_x\delta W=\text{lower order},\quad W=(n,\beta,v),
\end{equation}
\begin{equation}\label{eq:pencil}
 A^0=\begin{pmatrix}1&0&0\\0&1&0\\0&0&h\end{pmatrix},\qquad
 A^1=\begin{pmatrix}0&0&n\\0&0&0\\P_n&P_\beta&0\end{pmatrix}.
\end{equation}
Here $P_n=[w(1+w)\rho-2q]/n$ at fixed $\beta$ and $P_\beta=-2\kap\beta n^2$. For a covector $(-\omega,k)$,
\begin{align}
 \det(kA^1-\omega A^0)&=-\omega[h\omega^2-nP_n k^2],\qquad
 \lambda=\omega/k=0,\ \pm c_{\rm eff},\nonumber\\
 c_{\rm eff}^2&=\frac{nP_n}{h}=\frac{w(1+w)\rho-2q}{(1+w)\rho-2q}.\label{eq:charpoly}
\end{align}
For $h\ne0$, $c_{\rm eff}^2>0$ the three distinct eigenvalues have a complete basis. One choice is $r_\pm=(n,0,\pm c_{\rm eff})^T$, $r_0=(-P_\beta,P_n,0)^T$. The contact/polarization mode has $\delta P=0$; acoustic modes have $\delta\beta=0$. Equivalently in $(\rho,b,v)$,
\begin{equation}
 A^0=\operatorname{diag}(1,1,h),\quad
 A^1=\begin{pmatrix}0&0&(1+w)\rho\\0&0&b\\w&-2\kap b&0\end{pmatrix}.
\end{equation}
Both pencils give Eq.~\eqref{eq:charpoly}. An independently advected entropy variable adds a contact mode; it does not change the acoustic pair at fixed entropy and polarization. Retaining ideal common-congruence spin second moments also adds contact modes and leaves the same acoustic pair. Holding $b$ fixed during compression instead of $b/n$ would violate the selected spin transport and give a different, unjustified derivative.

On a fluid-comoving chart the coordinate acoustic speeds are $\dd r/\dd t=\pm(A/B)c_{\rm eff}$, and the material speed is zero. In a local Eulerian inertial frame with fluid speed $v$, the speeds are $v$ and $(v\pm c_{\rm eff})/(1\pm vc_{\rm eff})$. These expressions establish local causal acoustic cones when $0<c_{\rm eff}^2\leq1$, not global stability of collapse or well-posedness of our constrained gauge and boundary scheme. A regular hyperbolic metric gauge has light characteristics; coupling to that gravitational sector cannot remove a nonreal material acoustic pair by algebraic torsion elimination. We do not claim to have derived the characteristics of an unclosed multidimensional correlation theory.

\subsection{The complete constant-$w$ domain and its endpoints}
Let $z=q/\rho\geq0$ and $0\leq w<1$, $\rho>0$. Define
\begin{equation}
 z_s=\tfrac12w(1+w),\quad z_h=\tfrac12(1+w),\quad
 c_{\rm eff}^2-1=\frac{w^2-1}{1+w-2z}.
\end{equation}
The state classification is
\begin{longtable}{p{.25\textwidth}p{.27\textwidth}p{.37\textwidth}}
\toprule Domain & Material characteristics & Interpretation\\\midrule\endhead
$0\leq z<z_s$, $w>0$ & $h>0$, $0<c_{\rm eff}^2\leq w\leq1$ & Candidate smooth causal radial patch; no complete boundary or gauge proof follows.\\
$z=z_s$ & $c_{\rm eff}^2=0$, coincident eigenvalues & Primitive system is defective, not strongly hyperbolic at the endpoint.\\
$z_s<z<z_h$ & $h>0$, $c_{\rm eff}^2<0$ & Nonreal acoustic pair; high-frequency growth.\\
$z=z_h$ & $A^0$ singular & Euler acceleration cannot generically be evolved; undivided compatibility is required.\\
$z>z_h$ & $h<0$, $c_{\rm eff}^2>1$ & Real acoustic pair but outside the metric light cone; a causal-fluid interpretation fails.\\\bottomrule
\end{longtable}
For ordinary dust $w=0$, every $0<z<1/2$ is in the nonreal domain. In a local acoustic eigenmode, $\delta\eps\propto\exp(ikx+|k|\sqrt{-c_{\rm eff}^2}\,\tau)$. There is no wavelength-independent Sobolev growth estimate for the continuum initial-value problem. Hydrodynamics has a microscopic cutoff, but no specified cutoff/completion is present here; a finite grid or viscosity cannot establish well-posedness of these continuum equations. Analytic special solutions or restricted initial data can exist despite this generic smooth-data obstruction.

At $P_n=0$ and $h\ne0$, all three eigenvalues of $(A^0)^{-1}A^1$ vanish. Its rank is two when $P_\beta\ne0$, and its square is nonzero while its cube vanishes. It is a nilpotent Jordan system with only one eigenvector, not three stationary waves. For the coherent spinless dust endpoint ($w=\beta=0$), rank is one and the square vanishes. For Model A with independent variance perturbations, $P_\chi=-n^2$ remains nonzero even at $\chi=0$, so its formal three-variable pencil has rank two there; the zero-variance boundary of the physical state space needs a one-sided perturbation interpretation. Restricting exactly to the GR submodel removes that variance degree of freedom. GR dust has a useful pre-caustic Lagrangian description, but the primitive acoustic system is still not strongly hyperbolic with an equal-regularity norm. A repeated real root alone is not a viability check.

At $h=0$ the undivided radial Euler equation requires $\Dr P=0$, but no longer determines $a_r$. For $w<1$, $nP_n=(w-1)(1+w)\rho\ne0$ there, so generic gradients fail this constraint. Even specially tuned data satisfying it need a separate differential-algebraic continuation and cannot justify division by $h$. In the $h<0$ domain no universal ghost theorem is inferred from the fluid principal symbol; the direct conclusion is superluminal radial sound for this EOS.

\paragraph{Stiff exception.} At $w=1$, numerator and denominator coincide, so $c_{\rm eff}^2=1$ wherever $h\ne0$, including the negative-$h$ branch. At $q=\rho$ the original time matrix is still singular; cancelling the ratio does not invert it. Conserved $q$ and $\rho$ then both scale as the inverse square of cell volume, so $q/\rho$ is constant along each shell. There is no compression-driven crossing from an initially positive-enthalpy state to a spin-dominated state in this fixed-$w=1$ model. Special negative-effective-energy initial branches, for example homogeneous open geometries, are not excluded by this observation; their physical admissibility and boundary data are a different problem. We do not extend the $w<1$ acausality conclusion to them.

\subsection{General ordinary EOS: precisely what could change}
Let $\rho=\rho(n,\varsigma)$, $p=n\rho_n-\rho$, with entropy $\varsigma$ advected. At fixed $(\varsigma,\beta)$ let $h_0=\rho+p>0$ and $c_0^2=n p_n/h_0$. Keeping the same ideal spin correction and transport gives
\begin{equation}\label{eq:generalEOS}
 c_{\rm eff}^2=\frac{c_0^2 h_0-2q}{h_0-2q},\qquad
 c_{\rm eff}^2-1=\frac{(c_0^2-1)h_0}{h_0-2q}.
\end{equation}
For a subluminal ordinary EOS $0\leq c_0^2<1$, the negative-enthalpy branch is again superluminal. A positive-enthalpy causal patch requires $2q<c_0^2h_0$ (strictly for a nondegenerate acoustic pair). A microscopic high-density EOS can change $h_0,c_0$, spin susceptibilities and the source itself; no such EOS is supplied by specifying minimal EC gravity. Spin-dependent internal energy or independent polarization thermodynamics requires more derivatives and equations than Eq.~\eqref{eq:generalEOS}. We do not choose them to force a stable bounce.

\subsection{Can a bounce be reached inside the approved domain?}
For $w<1$, $\Dt\ln(q/\rho)=-(1-w)\Theta$: collapse increases $z$ along a spinning shell. A flat homogeneous bounce at $q=\rho$ has $c_{\rm eff}^2=w+2$ and passes the $z_s,z_h$ surfaces if it starts on the low-spin branch. The homogeneous ODE contains no acoustic gradients, so this does not invalidate the existing FLRW reference computations.

At a spherical volume turning point, the vorticity-free Raychaudhuri equation gives
\begin{equation}
 \Dt\Theta=-2\sigma^2-\tfrac\kap2[(1+3w)\rho-4q]+\nabla_\mu a^\mu
 \quad(\Theta=0).
\end{equation}
In a geodesic patch $\nabla_\mu a^\mu=0$, a volume minimum requires
\begin{equation}
 q>\tfrac14(1+3w)\rho+\sigma^2/\kap.
\end{equation}
For $0\leq w<1$ this lies above the positive-enthalpy causal bound because
\begin{equation}
 \frac{1+3w}{4}\rho-\frac{w(1+w)}2\rho
 =\frac{(1-w)(1+2w)}4\rho>0.
\end{equation}
Thus this selected closure has no geodesic, torsion-driven volume minimum inside that causal branch; shear makes the necessary spin threshold larger. This is a conditional inequality, not a theorem forbidding all spherical areal bounces. Inhomogeneous pressure gives $\nabla_\mu a^\mu=\Dr a_r+2\Gamma a_r/R+a_r^2$, which can alter a volume minimum, and $U=0$ is a distinct areal turning condition. A causal pressure-driven bounce, a shell-specific turn, or another microscopic EC source is not ruled out.

\section{Shell bounces, volume bounces and local regularity}
\subsection{Exact necessary shell condition for both selected models}
At a noncentral material shell define $U=\Dt R$. A candidate areal bounce needs $R>0$, $U=0$, and
\begin{equation}\label{eq:shellbounce}
 \Dt U=\Gamma a_r-\frac{M}{R^2}-\frac\kap2 PR>0.
\end{equation}
For either selected isotropic closure, wherever $h\ne0$ this is
\begin{equation}
 -\frac{\Gamma}{h}(w\Dr\rho-\Dr q)-\frac{M}{R^2}
 -\frac\kap2(w\rho-q)R>0.
\end{equation}
The spin term changes both the local pressure work and the constrained gravitational mass integral; $M$ cannot be replaced by an arbitrary local density. Spin-pressure gradients also affect acceleration. Ordinary pressure gradients can support a turn even without torsion. No independent spin bounce force is added. For a broader specified anisotropic source replace $P$ here by $P_r$ and use the full momentum equation including $\Dt Q$ and $P_r-P_\perp$; those terms cannot be chosen without matter closure.

Spatial curvature enters through $\Gamma^2=1+U^2-2M/R$ and the Hamiltonian constraint, not as an additional freely prescribed force. At $U=0$, $\Gamma^2=1-2M/R\geq0$. If $\Gamma\ne0$, the two null expansions have opposite signs, $\theta_\pm=\pm2\Gamma/R$; that turning sphere is not future trapped. The marginal $\Gamma=0$ case requires a regular-neck/chart analysis. This local observation neither identifies an event horizon nor excludes other portions of the evolution being trapped.

\subsection{Volume turning points and shear}
The local material cell volume is proportional to $BR^2$, not $R^3$:
\begin{equation}
 \Dt\ln(BR^2)=\Theta,\quad a_V\propto(BR^2)^{1/3},\quad
 \frac{\Dt^2a_V}{a_V}=\frac13\Dt\Theta+\frac19\Theta^2.
\end{equation}
At a shell areal extremum $H_\perp=0$, but $H_r$ need not vanish, so neither $\Theta=0$ nor a volume minimum follows. Conversely $\Theta=0$ need not imply $U=0$. The center $R=0$ is a regular limiting worldline, not a noncentral shell satisfying the stated $R>0$ bounce test.

The Levi--Civita vorticity-free Raychaudhuri equation follows from the kinematic Ricci identity and $R_{\mu\nu}u^\mu u^\nu=\kap(\eps+P_r+2P_\perp)/2$:
\begin{equation}
 \Dt\Theta=-\frac{\Theta^2}{3}-2\sigma^2-\frac\kap2(\eps+P_r+2P_\perp)
 +\nabla_\mu a^\mu.
\end{equation}
At a volume turning point for either selected model, a minimum requires
\begin{equation}\label{eq:volumebounce}
 \kap[4q-(1+3w)\rho]>4\sigma^2-2\nabla_\mu a^\mu,\quad
 \nabla_\mu a^\mu=\Dr a_r+\frac{2\Gamma a_r}{R}+a_r^2.
\end{equation}
Shear opposes the volume minimum. Pressure/spin gradients influence the acceleration divergence with no universal sign. The geodesic inequality in the characteristic section is only a restricted comparison; geodesic motion or zero shear is not imposed on the spherical field equations.

In FLRW, acceleration and shear vanish and all shells share one scale factor. At $H=0$ the constraint is $\kap(\rho-q)=3k/a^2$, and the volume/areal minimum additionally needs $4q>(1+3w)\rho$. Only for $k=0$ does the turning constraint give $q=\rho$. These are the appropriate averaged homogeneous reference conditions, not a coherent strong-symmetry FLRW solution.

\subsection{Necessary regularity is more than an inequality}
For a noncentral event with $R>0$, bounded $\rho,q,P$ keep the Ricci invariants finite:
\begin{align}
 \mathcal R&=\kap(\eps-3P)=\kap[(1-3w)\rho+2q],\\
 R_{\mu\nu}R^{\mu\nu}&=\kap^2(\eps^2+3P^2).
\end{align}
For a spherical isotropic no-flux source the Weyl invariant is
\begin{equation}
 C_{\mu\nu\rho\sigma}C^{\mu\nu\rho\sigma}
 =48\left(\frac{M}{R^3}-\frac{\kap\eps}{6}\right)^2,
\end{equation}
and the metric Kretschmann invariant is $C^2+2R_{\mu\nu}R^{\mu\nu}-\mathcal R^2/3$. These follow from the angular warped Riemann block and four-dimensional Weyl decomposition. The independent tests check the Weyl subtraction against the coordinate curvature. Bounded $M/R^3$ is therefore also needed; finite ordinary density alone does not bound all curvature. At a smooth center use regular limits and $M=O(R^3)$, not direct $0/0$ evaluation.

Model B additionally needs bounded spin and its covariant derivatives: the full curvature is
\begin{equation}
 \widetilde R^\rho{}_{\sigma\mu\nu}
 =R^\rho{}_{\sigma\mu\nu}
 +2\nabla_{[\mu}K^\rho{}_{\nu]\sigma}
 +2K^\rho{}_{[\mu|\lambda|}K^\lambda{}_{\nu]\sigma}.
\end{equation}
Model A's scalar variance alone does not fix full microscopic curvature statistics; bounded variance is not proof of every realization's regularity. For either model one must also maintain a nondegenerate material Jacobian, admissible characteristic domain, regular gauge and boundary, and actual local continuation. Finite samples of invariants or either inequality alone are not sufficient for nonsingular collapse, geodesic completeness, repeated oscillations or a complete exterior spacetime. No evolution or plots are produced in this theoretical stage.

\section{Initial data, regular center and exterior consistency}
\subsection{Independent initial data and physical restrictions}
After a radial relabelling $R_i=r$ on a monotone patch, independent profiles for the ideal constant-$w$ coherent branch are $U_i(r)$, $\rho_i(r)$ and signed $b_i(r)$. With fixed isentropic $C$, $n_i=(\rho_i/C)^{1/(1+w)}$ and $\beta_i=b_i/n_i$ are derived alternatives, not extra profiles. Entropy-dependent $C$ adds a declared advected entropy profile and its EOS. A microscopic finite spin per particle would bound $|\beta|$ by a material constant; neither its numerical value nor a polarization preparation is selected here.

Compute $\eps_i=\rho_i-\kap b_i^2$, $P_i=w\rho_i-\kap b_i^2$, then solve
\begin{align}
 M_i(r)&=\tfrac\kap2\int_0^r\eps_i(s)s^2\,\dd s,&
 \Gamma_i&=+\sqrt{1+U_i^2-2M_i/r},\\
 B_i&=1/\Gamma_i,&H_{r,i}&=U_i',\quad H_{\perp,i}=U_i/r,\\
 \Phi_i'&=-P_i'/h_i.
\end{align}
Set $M_i(0)=0$ and one lapse normalization, such as $\Phi(t,0)=0$. Require positive $B_i$, real $\Gamma_i$, smooth center limits, the chosen EOS and a declared characteristic domain. Collapse means $U_i<0$ for the desired noncentral shells; an initially trapped sphere is characterized by both null-expansion signs, not by the sign of $U$ alone. Do not independently prescribe mass, radial metric or extrinsic curvature after these constraints have fixed them. No manufactured spin profile is supplied as physical initial data.

At each later time, a smooth regular center requires
\begin{align}
 R&=\ell(t)r+O(r^3),& B&=\ell(t)+O(r^2),&
 U&=u_1(t)r+O(r^3),\\
 \rho&=\rho_c(t)+\rho_2(t)r^2+O(r^4),&
 b&=b_1(t)r+O(r^3),&
 \Phi&=\Phi_c(t)+\Phi_2(t)r^2+O(r^4).
\end{align}
Consequently $q=O(r^2)$, $M=\kap\rho_c\ell^3r^3/6+O(r^5)$, $\Gamma=1+O(r^2)$ and $\sigma^2=O(r^4)$. A coherent radial fluid can be regular, but cannot have nonzero coherent radial spin density exactly at the center. In ordinary dust the neighboring shells with nonzero $b$ still have the acoustic obstruction even if the center itself has $q=0$. Nor does $q_c=0$ prove absence of central volume reversal: the finite acceleration divergence can contain spin-pressure gradients. For example, at $R_i=r$, $a_r=2\kap b_1^2r/\rho_c+O(r^3)$ in the dust branch, yielding $\nabla_\mu a^\mu=6\kap b_1^2/\rho_c+O(r^2)$ at the center. This formal coefficient is not a viable bounce construction on a nonhyperbolic neighboring domain.

Particle and spin conservation imply $nBR^2$ and $bBR^2$ are material invariants. A finite-particle Jacobian must be tracked; a vanishing $R'$ alone is not shell crossing. At a regular closed-FLRW equator $R=a\sin\chi$, $B=a$, one has $R'=0$ without divergent matter or curvature. The divided monotone chart fails there and must not classify the event as a singularity.

\subsection{Analytical constraint-satisfying patches for both models}
As an analytical check, not a proposed astrophysical preparation, choose on $0\leq r<L$
\begin{equation}
 R_i=r,\quad U_i=-H_0r,\quad\rho_i=\rho_0>0,\quad q_i=q_0+q_2r^2,\quad H_0>0.
\end{equation}
Model A permits $q_0,q_2\geq0$. Model B requires $q_0=0$, $q_2=\kap b_1^2$ and signed $b_i=b_1r$. No new transport law is selected: these are only initial profiles subject to the already specified matter equations. The constraints integrate exactly:
\begin{align}
 M_i&=\kap\left[(\rho_0-q_0)r^3/6-q_2r^5/10\right],\\
 \Gamma_i^2&=1+\left[H_0^2-\kap(\rho_0-q_0)/3\right]r^2+\kap q_2r^4/5,\\
 B_i&=\Gamma_i^{-1},\qquad H_{r,i}=H_{\perp,i}=-H_0,\\
 h_i&=h_c-2q_2r^2,\quad h_c=(1+w)\rho_0-2q_0,\qquad
 \Phi_i=-\tfrac12\ln(h_i/h_c)
\end{align}
for a central lapse normalized to one and $h_c,h_i>0$. The last expression follows directly from $\Phi_i'=2q_2r/h_i$. These data are center regular. For $0<w<1$ a sufficiently small patch with $q_0+q_2L^2<w(1+w)\rho_0/2$ and $\Gamma_i^2>0$ also passes the local material characteristic test. Its initial zero shear is a feature of this example, not an assumption in the general equations, and it need not remain zero.

This is a constrained interior patch, not a matched finite star, proof of an evolution theorem or bounce. For example, a positive-density surface with $P_\Sigma=0$ would require $q_\Sigma=w\rho_\Sigma$, outside the nondegenerate causal low-spin domain for $0<w<1$. More generally at any such surface
\begin{equation}
 h_\Sigma=(1-w)\rho_\Sigma,\qquad
 c_{\rm eff,\Sigma}^2=-w\quad(0<w<1,\ \rho_\Sigma>0).
\end{equation}
This algebraic boundary obstruction applies to both selected closures even before the stronger coherent-contortion condition below. At $w=1$ the same cancellation instead gives $h=0$; at $w=0$ it forces $q=0$ and the pressureless endpoint.

\subsection{Metric traction and the vacuum exterior}
A source-free exterior in minimal EC has zero torsion by Cartan's algebraic equation. A locally spherical vacuum region is then a GR Schwarzschild region. For a timelike comoving material surface $\Sigma$, a source-layer-free effective metric junction has continuous induced metric and Levi--Civita extrinsic curvature. Projecting the Einstein equation onto the normal gives the necessary conditions
\begin{equation}
 P_\Sigma=0,\qquad M_\Sigma=M_{\rm ext},\qquad
 \Dt M_\Sigma=-\tfrac\kap2P_\Sigma R_\Sigma^2 U_\Sigma=0.
\end{equation}
These are necessary metric conditions, not the whole EC matching construction. In ordinary dust $P_\Sigma=-\kap b_\Sigma^2$, so metric vacuum matching itself requires $b_\Sigma=0$. With $w>0$ a cancellation $w\rho_\Sigma=\kap b_\Sigma^2$ can make metric pressure zero but leaves nonzero torsion; it is not a regular full-connection vacuum matching.

\subsection{A stronger boundary result derived from the full connection}
We can sharpen Stage 1's sufficient spin-free-collar discussion without claiming a general spinful shell formalism. Assume a non-null comoving boundary, continuous induced metric, a Darmois metric matching with no distributional Levi--Civita curvature, piecewise smooth bounded contortion, no singular spin sheet, and demand that the \emph{full} Riemann tensor have no Dirac delta. In Gaussian normal distance $\ell$ with outward spacelike normal $N_\mu$, the singular coefficient due to contortion alone is
\begin{equation}
 (\widetilde R^\rho{}_{\sigma\mu\nu})_{\rm sing}
 =\big(N_\mu[K^\rho{}_{\nu\sigma}]-N_\nu[K^\rho{}_{\mu\sigma}]\big)\delta(\ell),
 \qquad [K]=K_{\rm out}-K_{\rm in}.\label{eq:singular}
\end{equation}
This follows by differentiating a step in $\widetilde\Gamma=\Gamma+K$; bounded quadratic connection products add no delta. Work at a boundary point in an orthonormal frame with $u=e_{\hat0}$, $N=e_{\hat1}$ and $b_{\hat2\hat3}=b_\Sigma$. Our source dictionary gives $\Tr^{\hat0}{}_{\hat2\hat3}=-2\kap b_\Sigma$ and
\begin{equation}
 K^{\hat0}{}_{\hat2\hat3}=-\kap b_\Sigma.
\end{equation}
For a torsion-free exterior, Eq.~\eqref{eq:singular} has
\begin{equation}
 (\widetilde R^{\hat0}{}_{\hat3\hat1\hat2})_{\rm sing}
 =\kap b_\Sigma\delta(\ell).
\end{equation}
Hence $b_\Sigma=0$ is necessary for this stronger curvature-regular boundary. It is also sufficient to remove this contortion-jump delta within the stated radial branch, since all contortion components are proportional to $b$. Bounded one-sided derivatives and a compatible metric solution are still required. A whole spin-free collar is sufficient but not mathematically necessary for eliminating this delta; continuous $b\to0$ at the surface can suffice. Arbitrarily tapering $b$ during evolution would violate Eq.~\eqref{eq:radtransport}; any vanishing boundary value must be part of the initial-boundary problem and transported.

Importantly, direct contraction of this radial contortion-jump coefficient gives \emph{zero singular Ricci tensor}, even when its full Riemann coefficient is nonzero. Thus a source-layer-free effective Einstein equation alone need not detect this full-curvature singularity. This is exactly why we state the stronger boundary regularity assumption, rather than presenting $b_\Sigma=0$ as a universal theorem for all distributional EC matchings. This distinction agrees with the finite-curvature junction discussion in Sec. V of \cite{LC}; their static stellar solutions do not provide our dynamical existence result. A general interface with metric/contortion distributional cancellations or surface spin/stress remains a separate model. A delta spin sheet would make naive spin-squared products undefined without a regularized theory.

\subsection{Model A interface statistics and a remaining compact-edge obstruction}
For Model A, zero mean microscopic spin does not prove a regular spin-free junction: a finite variance or unresolved derivative/correlation layer may remain at the surface. Effective metric Darmois conditions and $P_\Sigma=0$ are necessary. A smooth microscopic spin-free collar satisfying the fluid equations is a sufficient route to ordinary GR matching, but this stage constructs no such collar. Vanishing nonnegative variance implies vanishing spin in the mean-square sense; it does not by itself bound microscopic derivatives or provide a full interface correlation model. Do not apply the coherent jump coefficient to an ensemble merely by replacing $b$ with $\sqrt{\langle b^2\rangle}$. Both models' material spin composition, when defined, is advected; a genuinely spin-free material outer region stays spin-free during smooth ideal evolution, but its pressure/EOS and metric matching still need to be solved.

For the stronger vacuum boundary with $w>0$, $b_\Sigma=0$ and $P_\Sigma=0$ require $\rho_\Sigma=0$. In an exact spin-free collar, Euler integrates to
\begin{equation}
 \Phi=-\frac{w}{1+w}\ln\rho+f(t).
\end{equation}
Thus a finite comoving zero-density edge is not automatically a regular finite-lapse boundary. If instead $\rho=Cn^{1+w}$ with bounded $\beta$ and $0<w<1$, then $q/\rho=(\kap\beta^2/C)n^{1-w}\to0$ at a vacuum edge. For regular polarization/entropy asymptotics the same leading logarithmic ordinary-fluid obstruction persists; a rapid singular composition variation would require its own analysis, not cure it. At $w=1$ a nondegenerate effective stiff coefficient has the analogous logarithmic issue. We have not constructed a finite isolated spinning star satisfying all these conditions. Possible alternatives are a microscopic EOS with a suitable surface, a specified non-vacuum exterior, or a derived interface model; none is selected ad hoc.

A Vaidya exterior contains null radiation. The selected radial coherent ideal interior has $Q=0$, so a comoving shell-free surface cannot be assigned an arbitrary radiating exterior merely because its interior mass has pressure work. Radiation needs a justified flux, noncomoving matter transfer or a specified layer. Full spinful junction laws, the surface action and global exterior causal completion are unresolved. A finite interior is therefore not a complete black-hole or baby-universe spacetime.

\section{Physical consistency, literature comparison and readiness}
\subsection{What the literature does and does not establish here}
\paragraph{Ideal Weyssenhoff.} Trautman \cite{Trautman} supplies the minimal variational setting; the specific ideal effective source/transport is selected from \cite{BHL,OK}. The latter is a constitutive field formulation, not a proof of a causal radial collapse problem for any EOS. Our spherical cancellation, complete reduced characteristic pencil, endpoint Jordan analysis, general-EOS derivative and contortion-jump test are derived consequences checked independently. They are not claimed as novel physics or a new fundamental no-bounce theorem.

\paragraph{Published collapse comparisons.} Hashemi et al. \cite{Hashemi} underlie Model A's homogeneous averaged closure. Pop{\l}awski \cite{Pop} uses a spherical thermodynamic setting with additional choices, including particle production in the discussed scenarios. A synchronous comoving Tolman gauge restricts acceleration and hence pressure gradients; it does not reproduce our arbitrary $\Phi(t,r)$ matter problem. Luz--Mena--Ziaie \cite{LMZ} study an \emph{effective dust} condition: ordinary pressure cancels spin pressure. That is different from ordinary $p=0$, $P=-q$. Neither their conclusions nor a static Weyssenhoff stellar solution in \cite{LC} supplies the missing radial hyperbolicity or exterior for our EOS.

\paragraph{Geometric well-posedness is not matter closure.} The recent analysis of Luz--Mena \cite{LWP} proves a geometric EC Cauchy result in a generalized harmonic gauge with compatible, sufficiently regular matter supplied ab initio; the stated hypotheses include conservation constraints on those supplied sources. Its non-strictly hyperbolic Leray--Ohya/Gevrey analysis is not a proof that our EOS/spin evolution yields those regular sources or a strongly hyperbolic numerical matter system. It does not contradict the explicit nonreal acoustic roots in Eq.~\eqref{eq:charpoly}. We do not claim EC geometry lacks a local well-posedness theory.

\paragraph{Modern spin hydrodynamics.} Weickgenannt \cite{Weick} derives a kinetic, dissipative spin theory and obtains causal, linearly stable modes about homogeneous equilibrium under specified matching and transport-coefficient conditions. Xie et al. \cite{Xie} analyze another dissipative spin framework and show that necessary small/large-wave-number conditions can miss finite-wave-number instabilities. These are distinct constitutive systems, not conflicting theorems about every first-order spin theory. They demonstrate the need for spin thermodynamics, microscopic matching and transport data, not a license to add relaxation/viscosity to ideal Weyssenhoff matter. We neither import their coefficients into EC nor establish their validity in dense gravitational collapse.

\subsection{Readiness decision for each formulation}
\begin{longtable}{p{.26\textwidth}p{.29\textwidth}p{.34\textwidth}}
\toprule Formulation & What is established & Decision\\\midrule\endhead
Model A: averaged GSM & Conditional local action, source, variance transport and characteristic pencil derived. & Algebraically closed under explicit averaging assumptions; same dust/enthalpy obstruction and unresolved finite boundary. The FLRW reference stays unchanged.\\
GR dust/perfect fluid & Standard spinless limits and analytical benchmarks. & Reference equations retained; a selected gauge/center/boundary problem precedes any new prototype.\\
Model B: coherent radial ideal fluid, $0<w<1$, $z<z_s$ & Formal closed continuum system and real causal material eigenbasis away from endpoints. & Candidate local smooth patch only; full coupled initial-boundary formulation and microscopic EOS still require justification. No prototype authorized by all six gates.\\
Both: ordinary dust, $0<z<1/2$ & Formal closure but imaginary acoustic pair. & Generic smooth radial evolution blocked.\\
Both: $h=0$ or $z>z_h$, $w<1$ & Singular time matrix or superluminal radial branch. & No causal bounce-through-domain solver justified.\\
Stiff ideal branch $w=1$ & Luminal pair for $h\ne0$; ratio fixed along shells. & Separate physical/energy/boundary problem; no generic regular-bounce claim.\\
Common-congruence statistical second moments & Conditional transport closes extra advected moments; same acoustic pair. & Not a stabilization mechanism. Physical averaging assumptions must be justified.\\
General spin/velocity/geometry correlations & Symmetry permits additional spherical sources. & Missing constitutive moments and evolution; no closed principal symbol or solver.\\\bottomrule
\end{longtable}
The task's six implementation prerequisites have \emph{not} all been met for a finite high-density full-Weyssenhoff collapse problem. We therefore stop before a spherical EC numerical prototype. This is the requested scientific result, not an unfinished implementation. A favorable local acoustic patch alone is insufficient to meet the gauge, EOS, center and exterior gates.

\subsection{What a justified simulation would record}
For any future accepted formulation, record $R$ and neighboring material shell labels, $U$, $M$, $\Gamma$, $B$, $\Phi$, $\rho,p,b,q$, all equation residuals and shell proper time $\tau_r=\int A\,\dd t$. Areal extrema require $R>0$, $U=0$ and the proper sign of $\Dt U$; a volume minimum instead requires $\Theta=0$, $\Dt\Theta>0$. A regular bounce additionally needs nonsingular matter/torsion/curvature, a continuing well-posed solution, finite particle Jacobians and compatible boundary data. An isolated turning point is not that result.

Track $\sigma^2$, $\Theta$, material caustics, noncentral $R\to0$, metric and full-connection curvature, and separate numerical/chart failure from physical divergence. Recollapse and oscillations require resolved extrema sequences with refinement and valid equations throughout. Geometric diagnostics remain
\begin{equation}
 \theta_\pm=\frac2R(U\pm\Gamma),\quad
 \theta_+\theta_-=\frac4{R^2}(2M/R-1),\quad
 N_R=\ln(R/R_i),\quad N_V=\tfrac13\ln(BR^2/B_iR_i^2).
\end{equation}
Both negative expansions identify future trapped spheres, both positive identify anti-trapping. A horizon-regular chart is needed for an implemented exterior; $2M/R=1$ is not itself a curvature singularity. Define accelerated volume e-folds, if used, by integrating $\Theta/3$ only when $\Theta>0$ and $\Dt\Theta/3+\Theta^2/9>0$, using each shell's proper time. Center values use regular limits. Export metric and spin-junction residuals; until a matching construction exists the status is \emph{not computed}. None of these local diagnostics establishes an event horizon, eternal oscillations or a baby universe.

\subsection{Open physics and the Dirac research direction}
The precise open tasks are: independently reconstruct the constrained full matter action; identify a microscopic high-density EOS and polarization susceptibility; justify a coarse-graining scale and close any retained spin/velocity/metric moments; establish the coupled gauge and characteristic problem; derive regular material surfaces or a justified non-vacuum exterior; and specify any surface stress/spin action without undefined products of distributions. These are separate obligations; minimal EC gravity alone does not solve them.

A minimally coupled Einstein--Cartan--Dirac action is a justified \emph{next research comparison}, since it specifies spin current and spin dynamics microscopically instead of guessing a fluid completion. Its spin source is a totally antisymmetric axial current, generally not the convective $u^\lambda b_{\mu\nu}$ source here. Eliminating its algebraic torsion produces a local axial-current contact interaction and a nonlinear Dirac equation. That interaction is algebraic in the spinors and does not change the Dirac principal symbol $\gamma^\mu\xi_\mu$; its characteristic cone remains metric-null. This structural observation is not a derivation of all gravitational constraints or a nonlinear existence theorem in a chosen spherical gauge.

Spherical spinors require careful symmetry of bilinears and sources; a single classical spinor is not a rotationally invariant scalar field. Torsion-free Einstein--Dirac work such as \cite{FSY} uses paired spin states to achieve spherical stress. In an EC extension, the \emph{total} axial current and cross-species interactions must be computed before imposing a singlet/multiplet ansatz. If it cancels identically for classical fields, the corresponding classical torsion effect cancels too. An unpolarized quantum state with zero mean current but nonzero current correlations is not equivalent to setting the classical current to zero before eliminating torsion. Quantum expectation values then require a controlled state, renormalization and correlation closure. Thus EC--Dirac offers a more definite starting action, not a guaranteed bounce or an automatic shortcut around the statistical-matter problem. No Dirac solver or invented four-fermion coefficient is implemented in this stage.

\subsection{Automated validation and its limits}
The new tests check the full reduced characteristic pencils, eigenvectors and endpoint Jordan structure, EOS derivative, sign domains and stiff exception; a nonspherical shear/acceleration control; direct radial divergence cancellation; spin transport and conserved particle polarization; center-series mass and lapse constraints; the full contortion-jump Riemann coefficient and its vanishing Ricci contraction; mixed Einstein components and the conditional geodesic bounce bound. Stage 1's 23 checks are rerun for geometry, Bianchi, mass/ADM constraints, conservation, Frenkel/source coefficients, FLRW, Schwarzschild and LTB limits, as are the existing application tests.

These are algebraic, analytical and numerical-evaluation validations, not a PDE time evolution, proof of microscopic validity or global well-posedness. Exact test inventory and environment versions are recorded with reproduction commands. Stage 1 files and the FLRW runtime remain intact; this directory is a new research stage.

\begin{thebibliography}{20}
\bibitem{audit} EC Collapse Laboratory, Stage 1 derivation and audit, repository commit \href{https://github.com/pdamkier-del/skovlab/tree/3efc96159537e834313ded3ba359624abaa8b26f/research/ec-spherical}{\texttt{3efc961}}, 10 November 2026.
\bibitem{Trautman} A. Trautman, Einstein--Cartan theory, \href{https://arxiv.org/abs/gr-qc/0606062}{gr-qc/0606062} (2006).
\bibitem{BHL} S. D. Br\'echet, M. P. Hobson and A. N. Lasenby, Weyssenhoff fluid dynamics in general relativity using a 1+3 covariant approach, \href{https://arxiv.org/abs/0706.2367v2}{0706.2367v2} (2007), Eqs. (19)--(28), (60)--(63), Appendix B.
\bibitem{OK} Yu. N. Obukhov and V. A. Korotky, The Weyssenhoff fluid in Einstein--Cartan theory, Class. Quantum Grav. 4 (1987) 1633, \href{https://doi.org/10.1088/0264-9381/4/6/021}{doi:10.1088/0264-9381/4/6/021}; original full action not independently reconstructed here.
\bibitem{Hashemi} M. Hashemi, S. Jalalzadeh and A. H. Ziaie, Collapse and dispersal of a homogeneous spin fluid in Einstein--Cartan theory, \href{https://arxiv.org/html/1407.4103v2}{1407.4103v2} (2015).
\bibitem{BB} C. G. B\"ohmer and P. Bronowski, The homogeneous and isotropic Weyssenhoff fluid, \href{https://arxiv.org/abs/gr-qc/0601089v2}{gr-qc/0601089v2} (2006).
\bibitem{Pop} N. Pop\l awski, Gravitational collapse of a fluid with torsion into a universe in a black hole, \href{https://arxiv.org/html/2008.02136v2}{2008.02136v2}, revision 20 March 2026.
\bibitem{LMZ} P. Luz, F. C. Mena and A. H. Ziaie, Influence of intrinsic spin in the formation of singularities for inhomogeneous effective dust space-times, \href{https://arxiv.org/abs/1811.12292v1}{1811.12292v1} (2019).
\bibitem{LC} P. Luz and S. Carloni, Static compact objects in Einstein--Cartan theory, \href{https://arxiv.org/abs/1907.11489}{1907.11489} (2019), Sec. V, especially the distinction between absence of a source layer and absence of distributional full curvature.
\bibitem{AKP} W. Arkuszewski, W. Kopczy\'nski and V. N. Ponomariev, Matching conditions in the Einstein--Cartan theory of gravitation, Commun. Math. Phys. 45 (1975) 183, \href{https://doi.org/10.1007/BF01629248}{doi:10.1007/BF01629248}; full component formalism not reproduced.
\bibitem{LWP} P. Luz and F. C. Mena, Local well-posedness of the initial value problem in Einstein--Cartan theory, \href{https://arxiv.org/abs/2505.22590v1}{2505.22590v1} (2025), introduction and Theorem 1; geometric theorem with matter supplied ab initio.
\bibitem{Weick} N. Weickgenannt, Linearly stable and causal relativistic first-order spin hydrodynamics, \href{https://arxiv.org/abs/2307.13561v2}{2307.13561v2} (2023).
\bibitem{Xie} X.-Q. Xie, D.-L. Wang, C. Yang and S. Pu, Causality and stability analysis for the minimal causal spin hydrodynamics, \href{https://arxiv.org/abs/2306.13880v3}{2306.13880v3} (2023).
\bibitem{FSY} F. Finster, J. Smoller and S.-T. Yau, Particlelike solutions of the Einstein--Dirac equations, \href{https://arxiv.org/abs/gr-qc/9801079}{gr-qc/9801079}, Phys. Rev. D 59 (1999) 104020; torsion-free static comparison, not a solution of the proposed EC collapse problem.
\bibitem{Brown} J. D. Brown, Action functionals for relativistic perfect fluids, \href{https://arxiv.org/abs/gr-qc/9304026}{gr-qc/9304026} (1993); effective conserved-current action, not a microscopic spin action.
\end{thebibliography}
\end{document}
